\documentclass[10pt,a4paper]{amsart}
\usepackage[margin=19mm]{geometry}
\usepackage{amsmath,amssymb,mathtools}
\usepackage[scr=rsfs,cal=euler]{mathalfa}
\usepackage{xfrac}
\usepackage[unicode,hidelinks,bookmarks=false]{hyperref}
\newcommand{\ie}{i.e.\ }
\newcommand{\cf}{cf.\ }
\newcommand{\resp}{respectively\ }
\newcommand{\Subsection}{\subsection}
\providecommand{\nobreakdash}{\nobreak\hbox{-}}
\setlength{\emergencystretch}{2em}
\multlinegap0pt
\usepackage{amssymb}
\usepackage{multirow}
\usepackage{algorithm}\usepackage{algpseudocode}
\usepackage{xcolor}
\usepackage[normalem]{ulem}
\newtheorem{assumption}{}
\newtheorem{lem}{}
\newtheorem{prop}{}
\newcommand{\jg}[1]{#1}
\title{Multigrid Monte Carlo Revisited: Theory and Bayesian Inference}
\author{Yoshihito Kazashi, Eike H. M\"uller, Robert Scheichl}
\date{Source: arXiv:2407.12149v3 (CC BY 4.0)}
\begin{document}
\maketitle

\noindent\emph{Source and licence.} Multigrid Monte Carlo Revisited: Theory and Bayesian Inference. Version arXiv:2407.12149v3 (CC BY 4.0), distributed by arXiv under the Creative Commons Attribution 4.0 International licence, http://creativecommons.org/licenses/by/4.0/, as recorded in the arXiv licence metadata for this exact version and shown on the abstract page https://arxiv.org/abs/2407.12149v3. Full licence notice: \texttt{licenses/arxiv-2407.12149-CC-BY-4.0.txt}.\par\smallskip

\noindent\textit{Editorial scope.} Counted unit: the article's prop prop:approximation-property (source0.tex lines 1467--1473). The article's complete original proof of it is delivered with it (source0.tex lines 2537--2597, one \texttt{proof} environment of 61 lines, which the article places away from the statement). No mathematics is written, repaired or re-derived: the statement and the proof are the article's own lines, in the article's own notation, and no retained line is substituted or re-flowed.  Retained uncounted as the source-local context the proof needs: lines 307--315; lines 376--378; lines 381--396; lines 1351--1353; lines 1355--1357; lines 1371--1384; lines 1374--1376; lines 1380--1382; lines 1401--1412; lines 1446--1452.  Nothing was omitted from any retained span, so the package carries no omission record.  No retained line contains a citation, so the wrapper carries no bibliography and no bibliography entry.  No retained line contains a figure environment, an image-inclusion command or drawing code, so no image is delivered and the package declares no asset dependency.  Editorial formatting only, in addition: an AMS \texttt{amsart} preamble that restates the article's own declarations and macros used by the retained lines, namely the environments \texttt{assumption}, \texttt{lem}, \texttt{prop} on separate counters, together with the standard packages those lines need.  The retained environments print in this document as Assumption 1, Lemma 1, Proposition 1 in order of appearance, whereas the article prints the same items with the numbering of its own full document.  The complete native source of the article as distributed in the arXiv e-print for this exact version is bundled unchanged as \texttt{sources/162432/source0.tex}.
\begin{assumption}
	\label{assu:Pell}There is a decreasing function
        $\Phi\colon\mathbb{N}_{0}\to[0,\infty)$ such that %
        $P_{\ell}\colon\mathbb{R}^{n_{\ell}}\to V_{\ell}$ satisfies
	\begin{equation}
		c_{1}\|P_{\ell}x_{\ell}\|_{H}\leq\Phi(\ell)\|x_{\ell}\|_{2}\leq c_{2}\|P_{\ell}x_{\ell}\|_{H}\quad\text{for all }x_{\ell}\in\mathbb{R}^{n_{\ell}},\label{eq:Pl-property}
	\end{equation}
	with some constants $c_1,c_2>0$ independent of $\ell$.
\end{assumption}

\begin{equation}
	a(u,\varphi)=\langle f,\varphi\rangle_{H}\text{ for all }\varphi\in V,\label{eq:eq-given-by-A}
\end{equation}

\begin{assumption}
	\label{assu:WtoH-best-Vell-approx}There exists a subspace 
	$W\subset V \subset H$ equipped with a norm $\|\cdot\|_W$ such that
        \begin{enumerate}
        \item[(a)] the solution $u$ of \eqref{eq:eq-given-by-A}
	belongs to $W$ and $\|u\|_{W}\leq C_{\mathcal{A}}\|f\|_{H}$,
	for some constant $C_{\mathcal{A}}>0$;
        \item[(b)] the best-approximation error in $V_\ell\subset V$ with respect to
        the $V$-norm satisfies
	\[
        \inf_{w_{\ell}\in  V_{\ell}}\|w-w_{\ell}\|_{V}\leq\sqrt{\Psi(\ell)}\|w\|_{W}\,,\qquad\text{for
                  any }w\in W 
        \]
        for a  decreasing function $\Psi\colon\mathbb{N}_{0}\to[0,\infty)$ \jg{with $\Psi(\ell)\rightarrow 0$ as $\ell\rightarrow \infty$}.
              \end{enumerate}
\end{assumption}

\begin{equation}
	a(u,\varphi)+b(u,\varphi)=\langle f,\varphi\rangle_{H}=\text{ for all }\varphi\in V;\label{eq:eq-perturbed-original}
\end{equation}

\begin{equation}
	a(u_{\ell},\varphi_{\ell})+b(u_{\ell},\varphi_{\ell})=\langle f,\varphi_{\ell}\rangle_{H}\text{ for all }\varphi_{\ell}\in V_{\ell}.\label{eq:eq-perturbed-Vell}
\end{equation}

\begin{lem}\label{lem:approximation-perturbed}
	If $\mathcal{A}$ satisfies Assumption~\ref{assu:WtoH-best-Vell-approx},
	then so does the perturbed operator $\widetilde{\mathcal{A}}=\mathcal{A}+\mathcal{B}^{*}\Gamma^{-1}\mathcal{B}$ and
	\begin{equation}
		\|u\|_{W}  \leq C_{\widetilde{\mathcal{A}}}\,\|f\|_{H}\label{eq:perturbedu-Wnorm-bound}
	\end{equation}
	holds. Moreover, the solution $u\in V$ of (\ref{eq:eq-perturbed-original})
	can be approximated by the solution $u_{\ell}\in V_{\ell}$ of (\ref{eq:eq-perturbed-Vell})
	with an error
	\begin{equation}
		\|u-u_{\ell}\|_{H}\leq C\Psi(\ell)\|u\|_{W},\label{eq:perturbed-H-W-bound}
              \end{equation}
          where $\Psi(\ell)$ is defined in Assumption~\ref{assu:WtoH-best-Vell-approx}.
\end{lem}

	\begin{equation}
		\|u\|_{W}  \leq C_{\widetilde{\mathcal{A}}}\,\|f\|_{H}
	\end{equation}

	\begin{equation}
		\|u-u_{\ell}\|_{H}\leq C\Psi(\ell)\|u\|_{W},
              \end{equation}

\begin{assumption}\label{ass:inverse_inequality}
	There exists a constant $C_{\Psi}>1$, such that the function
	$\Psi\colon\mathbb{N}_{0}\to[0,\infty)$ in Assumption~\ref{assu:WtoH-best-Vell-approx} satisfies
	\begin{equation}
		\Psi(\ell-1)\leq C_{\Psi}\Psi(\ell) \qquad\text{for all $\ell\geq1$}
		\label{eqn:inverse_inequality1}
	\end{equation}
	and the following inverse inequality holds:
        \begin{equation}\label{eqn:inverse_inequality2}
          \|\varphi_{\ell}\|_{V}\leq\|\varphi_{\ell}\|_{H}/\sqrt{\Psi(\ell)}, \quad \text{for all $\varphi_{\ell}\in V_{\ell}\,$}.
          \end{equation}
\end{assumption}

\begin{prop}
	\label{prop:1/tildeAll-lb} Let Assumptions~\ref{assu:Pell} and \ref{ass:inverse_inequality} be satisfied. Then, there exists a constant $\tilde{C}>0$ such that 
	\begin{equation}
          \frac{\Psi(\ell-1)}{(\Phi(\ell))^{2}}\leq\frac{\tilde{C}}{\|\widetilde{A}_{\ell}\|_{2}}
          \qquad\text{for all $1\leq \ell\leq L$.}\label{eq:inv-norm-bound}
	\end{equation}
\end{prop}

\begin{prop}
	\label{prop:approximation-property}Let Assumptions~\ref{assu:Pell}, \ref{assu:WtoH-best-Vell-approx} and \ref{ass:inverse_inequality} hold. Then, we have the approximation property for the perturbed matrix $\widetilde{A}_{\ell}$
	\begin{align*}
		\|\widetilde{A}_{\ell}^{-1}-I_{\ell-1}^{\ell}\widetilde{A}_{\ell-1}^{-1}I_{\ell}^{\ell-1}\|_{2}
		 & \leq\frac{C}{\|\widetilde{A}_{\ell}\|_{2}}.
	\end{align*}
\end{prop}

\begin{proof}[Proof of Proposition~\ref{prop:approximation-property}]
	Let $f_{\ell}\in\mathbb{R}^{n_{\ell}}$ be given. For
	$\widetilde{A}_{\ell}=A_{\ell}+B_{\ell}\Gamma^{-1}B_{\ell}^{\top}$ and
	$\widetilde{A}_{\ell-1}=A_{\ell-1}+B_{\ell-1}\Gamma^{-1}B_{\ell-1}^{\top}$, let
	\begin{align*}
		u_{\ell} & :=\widetilde{A}_{\ell}^{-1}f_{\ell}\in\mathbb{R}^{n_{\ell}}
		\quad\text{and}\quad
		u_{\ell-1}:=\widetilde{A}_{\ell-1}^{-1}I_{\ell}^{\ell-1}f_{\ell}\in\mathbb{R}^{n_{\ell-1}}.
	\end{align*}
	With $u_{\ell}$ and $u_{\ell-1}$, let $\tilde{u}_{\ell}:=P_{\ell}u_{\ell}$
	and $\tilde{u}_{\ell-1}:=P_{\ell-1}u_{\ell-1}$, so that
	\begin{align}
		\|(\widetilde{A}_{\ell}^{-1}- & I_{\ell-1}^{\ell}\widetilde{A}_{\ell-1}^{-1}I_{\ell}^{\ell-1})f_{\ell}\|_{2}  =\|u_{\ell}-I_{\ell-1}^{\ell}u_{\ell-1}\|_{2} =\|P_{\ell}^{-1}\tilde{u}_{\ell}-I_{\ell-1}^{\ell}P_{\ell-1}^{-1}\tilde{u}_{\ell-1}\|_{2} \\
		                              & =\|P_{\ell}^{-1}(\tilde{u}_{\ell}-\tilde{u}_{\ell-1})\|_{2}
		\leq c(\|\tilde{u}_{\ell}-u\|_{H}+\|u-\tilde{u}_{\ell-1}\|_{H})/\Phi(\ell),\label{eq:diff-in-uell-uellminus1}
	\end{align}
	where we used Assumption~\ref{assu:Pell}.

	To use (\ref{eq:perturbed-H-W-bound}) we rewrite the equation $\widetilde{A}_{\ell}u_{\ell}=f_{\ell}$
	in $\mathbb{R}^{n_{\ell}}$ and $\widetilde{A}_{\ell-1}u_{\ell-1}=I_{\ell}^{\ell-1}f_{\ell}$
	in $\mathbb{R}^{n_{\ell-1}}$ in the variational problems in function
	spaces $V_{\ell}$ and $V_{\ell-1}$.
	Let $(P_{\ell}^{-1})^{*}$ be
	the adjoint of $P_{\ell}^{-1}\colon(V_{\ell},\langle\cdot,\cdot\rangle_{H})\to\mathbb{R}^{n_{\ell}}$,
	i.e., $\langle(P_{\ell}^{-1})^{*}x_{\ell},y_{\ell}\rangle_{H}=x_{\ell}^{\top}P_{\ell}^{-1}y_{\ell}$
	for all $x_{\ell}\in\mathbb{R}^{n_{\ell}}$, $y_{\ell}\in V_{\ell}$,
	and define $F_{\ell}(\varphi):=\langle(P_{\ell}^{-1})^{*}f_{\ell},\varphi\rangle_{H}$
	for $\varphi\in V$. Then, the function $\tilde{u}_{\ell}=P_{\ell}u_{\ell}$
	satisfies
	\begin{align*}
		a(z_{\ell},\tilde{u}_{\ell})+b(z_{\ell},\tilde{u}_{\ell}) & =f_{\ell}^{\top}P_{\ell}^{-1}z_{\ell}=\langle(P_{\ell}^{-1})^{*}f_{\ell},z_{\ell}\rangle_{H}  =F_{\ell}(z_{\ell})\qquad\text{ for all }z_{\ell}\in V_{\ell},
	\end{align*}
	and $\tilde{u}_{\ell-1}=P_{\ell-1}u_{\ell-1}$ satisfies
	\begin{align*}
		a(z_{\ell-1},\tilde{u}_{\ell-1})+b(z_{\ell-1},\tilde{u}_{\ell-1}) & =(I_{\ell}^{\ell-1}f_{\ell})^{\top}P_{\ell-1}^{-1}z_{\ell-1}=(f_{\ell})^{\top}I_{\ell-1}^{\ell}P_{\ell-1}^{-1}z_{\ell-1} \\
		                                                  & =\langle(P_{\ell}^{-1})^{*}f_{\ell},z_{\ell-1}\rangle_{H}                                                                          =F_{\ell}(z_{\ell-1})\text{ for all }z_{\ell-1}\in V_{\ell-1}.
	\end{align*}
	Hence, $\tilde{u}_{\ell}$ and $\tilde{u}_{\ell-1}$ are approximations of the solution
	$u$ to the problem
	\[
		a(u,\varphi)+b(u,\varphi)=F_{\ell}(\varphi)\quad\text{for all}\quad\varphi\in V.
	\]
	We now use \eqref{eq:perturbedu-Wnorm-bound} and \eqref{eq:perturbed-H-W-bound} in Lemma~\ref{lem:approximation-perturbed} to
	\eqref{eq:diff-in-uell-uellminus1} to obtain
	\begin{align*}
		\|(\widetilde{A}_{\ell}^{-1}-I_{\ell-1}^{\ell}\widetilde{A}_{\ell-1}^{-1}I_{\ell}^{\ell-1})f_{\ell}\|_{2} & \leq c(\Psi(\ell)+\Psi(\ell-1))\|u\|_{W}/\Phi(\ell)                        \leq c\Psi(\ell-1)\|(P_{\ell}^{-1})^{*}f_{\ell}\|_{H}/\Phi(\ell).
		%\\& \leq c\Psi(\ell-1)\|(P_{\ell}^{-1})^{*}f_{\ell}\|_{V}/\Phi(\ell).
	\end{align*}
	Finally, Assumption~\ref{assu:Pell} implies $\|P_{\ell}^{-1}\varphi_{\ell}\|_{2}\leq\frac{c_{2}}{\Phi(\ell)}\|\varphi_{\ell}\|_{H}$
	for $\varphi_{\ell}\in V_{\ell}$ and thus
	\[
		\|(P_{\ell}^{-1})^{*}\|_{\mathbb{R}^{n_{\ell}}\to V_{\ell}}=\|P_{\ell}^{-1}\|_{V_{\ell}\to\mathbb{R}^{n_{\ell}}}\leq\frac{c_{2}}{\Phi(\ell)},
	\]
	where $\|P_{\ell}^{-1}\|_{V_{\ell}\to\mathbb{R}^{n_{\ell}}}$ denotes
	the operator norm from $(V_{\ell},\langle\cdot,\cdot\rangle_{H})$
	to $\mathbb{R}^{n_{\ell}}$. We now invoke Proposition~\ref{prop:1/tildeAll-lb}
	to conclude
	\[
		\|(\widetilde{A}_{\ell}^{-1}-I_{\ell-1}^{\ell}\widetilde{A}_{\ell-1}^{-1}I_{\ell}^{\ell-1})\|_{2}\leq c\Psi(\ell-1)/(\Phi(\ell))^{2}\leq\frac{C}{\|\widetilde{A_{\ell}}\|_{2}}.
	\]
\end{proof}

\end{document}
